\section*{Related work}
Healthcare NLP has a long methodological base in converting unstructured biomedical and clinical text into analyzable evidence. General introductions and reviews emphasize both the promise of text mining and the fragility of systems that are evaluated only in narrow settings~\cite{nadkarni2011nlp,zweigenbaum2007frontiers,kreimeyer2017nlpsystems,wang2018clinicalie}. Restricted clinical-note resources and shared tasks such as MIMIC-III and i2b2/VA have shaped clinical NLP, but they require credentialed access or task-specific permissions and were therefore not used in this public CPU-only project~\cite{johnson2016mimic,uzuner2011i2b2}.

Public biomedical abstract corpora have made sentence classification a reproducible benchmark problem. PubMed-PICO was introduced for PICO element detection in medical text and has been used to study neural sentence classification with long short-term memory networks~\cite{jin2018pico}. PubMed 20k/200k RCT provides section-labelled randomized-trial abstracts for sequential sentence classification~\cite{dernoncourt2017pubmed}. EBM-NLP and related work show the broader value of patient, intervention, outcome, and evidence-bearing sentence annotations for systematic-review support~\cite{nye2018ebmnlp,kim2011ebm}. NICTA-PIBOSO is a manually annotated corpus of medical abstracts~\cite{amini2018nicta}, but it was not usable as an automatic external validation source in this run because the accessible archive did not contain usable files and the latest record was restricted. Other public biomedical tasks, including PubMedQA~\cite{jin2019pubmedqa} and adverse-event extraction corpora~\cite{gurulingappa2012ade}, were considered but were less directly aligned with a CPU-feasible calibrated binary sentence-classification benchmark.

The evidence-synthesis literature also motivates this task. Automated citation screening, review automation, and risk-of-bias systems show how NLP can reduce workload while still requiring transparent human oversight~\cite{cohen2006workload,tsafnat2014automation,marshall2014risk,marshall2015risk,marshall2015robotreviewer}. The present study is narrower than full systematic-review automation: it tests sentence-level rhetorical detection rather than eligibility, effect-size extraction, or risk-of-bias judgment.

Classical machine-learning methods remain strong baselines for text classification. TF-IDF term weighting and machine-learning text categorization have a long empirical history~\cite{salton1988term,sebastiani2002text}, and support vector machines provide a well-established margin-based classifier~\cite{cortes1995svm,joachims1998text}. Efficient linear solvers and probabilistic calibration methods make sparse models practical for large biomedical corpora~\cite{fan2008liblinear,platt1999probabilistic}. The classical models in this study were implemented with scikit-learn~\cite{pedregosa2011sklearn}.

Recurrent neural networks provide a contrasting sequence-aware modelling family. LSTM architectures were designed to address long-range dependency learning in sequential data~\cite{hochreiter1997lstm}, and bidirectional recurrent networks can incorporate context from both token directions~\cite{schuster1997brnn}. Static word vectors, convolutional sentence models, contextual embeddings, and Transformer language models have substantially shaped modern NLP~\cite{mikolov2013word2vec,kim2014cnn,peters2018elmo,vaswani2017attention,devlin2019bert}. Biomedical and scientific variants such as SciBERT, BioBERT, and PubMedBERT further demonstrate the value of domain-specific pretraining~\cite{beltagy2019scibert,lee2020biobert,gu2021pubmedbert}. This benchmark nevertheless focused on LSTM/BiLSTM neural baselines because they are realistic CPU-only comparators; GPU-scale Transformer fine-tuning was outside scope.

Evaluation in this setting must go beyond accuracy. The Brier score is a proper scoring rule for probabilistic forecasts~\cite{brier1950verification}, ROC and precision-recall analyses capture complementary aspects of binary discrimination~\cite{fawcett2006roc,davis2006prroc,saito2015pr}, and post-hoc calibration methods were developed to convert classifier scores into better probability estimates~\cite{niculescu2005probabilities,zadrozny2002classifier,guo2017calibration}. Clinical prediction literature has argued that calibration should be assessed alongside discrimination because useful ranking is not the same as reliable probabilities~\cite{vancalster2016hierarchy,vancalster2019calibration,austin2019ici}.

Transparent reporting and reproducibility are especially important when healthcare-adjacent machine learning is evaluated on proxy labels. TRIPOD and related guidance emphasize validation, complete reporting, and cautious interpretation~\cite{collins2015tripod,moons2015tripod,luo2016mlguidelines}. FAIR data principles, reproducibility recommendations, model cards, and datasheets provide complementary expectations for documenting code, data, and model limitations~\cite{wilkinson2016fair,pineau2021reproducibility,mitchell2019modelcards,gebru2021datasheets}. Clinical AI guidance, including MI-CLAIM, SPIRIT-AI, CONSORT-AI, and DECIDE-AI, reinforces that benchmark performance is not deployment evidence~\cite{norgeot2020miclaim,rivera2020spiritai,liu2020consortai,vasey2022decideai}. Broader healthcare AI commentary similarly warns that clinical impact requires validation, workflow integration, ethics review, and prospective evaluation~\cite{rajkomar2019medicine,kelly2019challenges,chen2017prediction,wiens2019roadmap,char2018ethical,nagendran2020clinicians}.
